% page 502 of the facsimile (image p-501 of the scan)
% block 157.5 x 246.3 mm ; left margin 8.0 mm
\pgc{-12.700mm}{0.000mm}{
\l{\cel{3}494}
\l{}
\l{\sou{\textquotedbl{}requiem\textquotedbl{}}. (\sou{Liturgio katolika}). I. Prego por la mortinti.}
\l{\cel{3}- II.\cel{1}\sou{Meso di \textquotedbl{}requiem\textquotedbl{}}\cel{1}:Meso por la repozo di la anmo}
\l{\cel{3}di mortinto. - III. Parti di la mortinto-meso, igita}
\l{\cel{3}muzikajo. - .}
\l{}
\l{\sou{requizitar}. (\sou{trans.}) (\sou{aludante la stat-autoritatozi.}) .}
\l{\cel{3}Postular ke persono o kozi esez igata en la dispono di}
\l{\cel{3}ta autoritatozi por servado od uzado publiko-destine}
\l{\cel{3}(+\sou{publico}-destine). - DEFIRS.}
\l{}
\l{\sou{resekar}. (\sou{trans.})\cel{2}(\sou{kirurgio}).Ablacionar un ek la extremji}
\l{\cel{3}di osto malada. -}
\l{}
\l{\sou{resipicenca}r (\sou{netrans.})\cel{2}Repentar ye maniero qua eventigas la}
\l{\cel{3}ri-komenco agar etiko-konforme. - DEFI.}
\l{}
\l{\sou{reskripto}. I. (yuro-cienco, che la Romani antiqua).Respondo,}
\l{\cel{3}da la imperiestro, a la questioni agita a lu da la guberni-}
\l{\cel{3}estri, da la judiciisti, e c., pri ula problemi solvenda.}
\l{\cel{3}- II. (en\cel{1}\sou{nia landi, nun)}. Dekreto da la stato-chefo.-}
\l{\cel{3}III. (\sou{Vatikano}).Averto papala, qua decidas pri ta o ca}
\l{\cel{3}yuro-punto. - DEFIRS.}
\l{}
\l{\sou{resonar}. (\sou{netrans}.)I. Retrosendar la soni per reflekto. -}
\l{\cel{3}II. Produktar soni qui igesas pluforta per la reflekto}
\l{\cel{3}interne di klozajo (+kluzajo) di qua la kapaceso ne esas}
\l{\cel{3}sat granda pri facar eko. - DEFIS.}
\l{}
\l{\sou{resorbar}. (\sou{trans.}) (\sou{medic}.)Igar retroirar aden la cirkulado}
\l{\cel{3}(liquido quan la vaskuli lasabis eskapar). - DEFIS.}
\l{}
\l{\sou{resorto}. - I. Peco di mekanismo qua, per sua destenso, mo-}
\l{\cel{3}vigas peco vicina. - II. (\sou{metaf}.)To quo impulsas. - FRS.}
\l{}
\l{\sou{resortisar}. (\sou{netrans., ad).}\cel{1}Esar judicienda da ta o ca tri-}
\l{\cel{3}bunalo. - DEFRS.}
\l{}
\l{\sou{respektar}. - (\sou{trans.}) I. Egardar kom kozo quan onu devas}
\l{\cel{3}atencar pro lua importo. - II. Egardar kom kozo quan onu}
\l{\cel{3}devas ne domajar, ne lezar, pro lua bonegeso. - III. Honor-}
\l{\cel{3}izar (ulu) per deferencar profunde lu, pro lua karaktero,}
\l{\cel{3}pro lua alta situeso en la hierarkio. - DEFIRS.}
\l{}
\l{\sou{respir}ar. (\sou{trans. e netrans.})\cel{2}Aspirar aero por regenerar la}
\l{\cel{3}sango, e, pose, expirar lu. - DEFIS.}
\l{}
\l{\sou{responda}r.(trans., ad).\cel{2}Sendar, direktar ad ulu di qua onu}
\l{\cel{3}recevis demando, questiono, akuzo, e c. to quon onu havas}
\l{\cel{3}dicenda kambie. - EFIS.}
\l{}
\l{\sou{responsa}r (\sou{netrans., pri})Obligesar reparor, per indemno, la}
\l{\cel{3}detrimento o domajo quin onu agis od agos a kunhomo. -EFIS.}
\l{}
\l{\sou{responsorio}. (\sou{liturgio}\cel{1}katol.)\cel{2}Citaji ek biblo, dicata o kan-}
\l{\cel{3}tata pos omna leciono di matutino. - DEFIS}
}
